\chapter{Introduction}

This chapter explains why we are proposing a smart bin system. It starts with how waste management in Kampala has reached its present state, moves to the problem that follows from this, and then gives the objectives, the scope and the significance of the study.

\section{Background}

Managing waste involves more than taking rubbish out of a bin. The waste has to be separated, stored, collected and taken to a suitable place, and when these steps are poorly connected, materials that could have been recovered end up mixed with everything else. The World Bank estimated that the world produced 2.01 billion tonnes of municipal solid waste in 2016 and expects 3.40 billion tonnes by 2050 (Kaza et al., 2018).

Kampala shows the problem at city level. KCCA puts the waste generated in the city at over 2,500 tonnes a day (KCCA, 2026). Mwesigwa (2024) reports that about 1,200 tonnes a day are collected by garbage trucks, which is roughly 48 percent, and that KCCA's own collection rate rose from 35 percent in 2016 to 59 percent in 2023. KCCA reported 65 percent for the 2023/2024 financial year (KCCA, 2024a). These figures are not the same because they are measured in different ways, but they all show that a large part of the waste is never collected. Waste that is left behind often ends up in drainage channels and on roads, where it blocks water and causes flooding (Ngwomoya, 2018; Mwesigwa, 2024).

The reasons for the gap have built up over several years. Collection is shared between KCCA and private contractors (KCCA, n.d.). Reports on the arrangement mention too few trucks and frequent breakdowns, limited funding, contractors who collect on irregular schedules or overcharge, and residents who do not pay collection fees (Ngwomoya, 2018; Mwesigwa, 2024). The city has also kept growing. KCCA lists population growth, economic growth and rising waste generation among the reasons its main landfill filled up (KCCA, 2026).

Most of the waste is also not separated at source. KCCA's own figures put the composition at about 85 percent organic waste, 8 percent plastic, 2 percent paper and 5 percent other material (KCCA, 2026). Material with a use, such as food waste that could go to farms or compost, is therefore likely to be carried to the same dump as everything else. KCCA has asked people to separate organic and inorganic waste at source (KCCA, n.d.). Okello (2024) names low public awareness, weak enforcement and a lack of reliable data among the problems, and Mwesigwa (2024) suggests that degradable waste could be used to produce biogas.

The Kiteezi landfill is where these problems ended up. It opened in 1996 on 39 acres and was receiving about 1,200 tonnes a day when it was closed (KCCA, 2026). According to Mwesigwa (2024), it had run for about ten years longer than it was meant to, and Africanews (2024) reports that plans to move waste disposal elsewhere had stalled since 2016. After heavy rain in August 2024 a large part of the waste mound collapsed onto the houses next to it (Petley, 2024; Okello, 2024). KCCA reported that at least 35 people died (KCCA, 2024b). A year later the site was no longer in official use but was still giving off methane, and disposal had moved to a new site in Mpigi that the National Forestry Authority said encroaches on a protected area (Agence France-Presse, 2025). Among the reforms Okello (2024) recommends are waste diversion strategies and more awareness of waste separation. We take from this that a lot of what reaches the dump could have been kept out of it if it had been separated and collected on time at the places where it is produced.

The Internet of Things (IoT) connects physical devices that take measurements to other systems that can use them. For waste, Medvedev et al. (2015) describe a decision support system that combines RFID tags, sensors, cameras and actuators so that truck drivers can share information in real time and routes can be changed as conditions change. In smart bins, fill-level sensors report how much space is left. Ahmed et al. (2022) used time series from IoT-connected municipal bins to forecast the status of bins with four deep learning models, and found that an LSTM model gave the lowest percentage error (MAPE of 1.855). Image classification has also been moved onto small devices. White et al. (2020) report an accuracy of 0.970 on the TrashNet test set with a model designed to run on a Jetson Nano. Li and Grammenos (2022) report 95.98 percent on a seven-class test set with a MobileNet V3 model on the same board, running at 40 inferences a second, and show a cheaper K210 board that uses far less power. These results come from their own datasets and conditions, so we cannot assume the same figures for hotels and restaurants in Uganda.

Some smart bin work has already been done locally and in the region. Musoke (2022) describes a prototype with fill-level sensors, material detection and reporting to a cloud platform, and Turyahebwa (2023) describes camera-based identification, sorting into compartments and remote reporting over LTE. Both are Makerere projects, and we have only seen the repository abstracts. They show that this kind of work exists locally, but not how well the prototypes performed or whether they were ever used outside the lab. We will read the full reports during the study of the current situation. Kwagala and Badoru (2023) propose a smart municipal waste system for Kampala that uses smart bins for billing and route optimisation, but the results in their abstract are expected results and were not yet measured. Komane and Mathonsi (2023) design a system for Johannesburg with ultrasonic sensors, GPS and GSM, and report no measured accuracy for it.

The image data also has to match the waste classes we plan to use. TrashNet has 2,527 images of glass, paper, cardboard, plastic, metal and general trash, but nothing labelled as organic (Yang and Thung, 2016; White et al., 2020). It may help with early experiments on some classes, but we will have to collect our own images, especially of food waste and of items photographed under the conditions the prototype will work in.

We chose hotels and restaurants because they prepare food and use a lot of packaging. The preliminary study will establish which types of waste the participating sites produce, how their staff separate it and how collection is arranged. We will also look at space, cleaning, electricity and network access, since a bin that is hard to use or keep clean may be ignored even if its electronics work.

Our proposed system brings classification, compartment monitoring, fill-up prediction and collection support together in one prototype. A mobile application will show the status of the bin, alert authorised collectors and record where the separated waste is taken. By smart disposal we mean sending each recognised item to the right compartment and arranging for it to be taken to a suitable recipient. We are not claiming that the bin treats the waste itself, or that every item can be recycled.
